%%% The main file. It contains definitions of basic parameters and includes all other parts.

%% Settings for single-side (simplex) printing
% Margins: left 40mm, right 25mm, top and bottom 25mm
% (but beware, LaTeX adds 1in implicitly)
\documentclass[12pt,a4paper]{report}
\setlength\textwidth{145mm}
\setlength\textheight{247mm}
\setlength\oddsidemargin{15mm}
\setlength\evensidemargin{15mm}
\setlength\topmargin{0mm}
\setlength\headsep{0mm}
\setlength\headheight{0mm}
% \openright makes the following text appear on a right-hand page
\let\openright=\clearpage

%% Settings for two-sided (duplex) printing
% \documentclass[12pt,a4paper,twoside,openright]{report}
% \setlength\textwidth{145mm}
% \setlength\textheight{247mm}
% \setlength\oddsidemargin{14.2mm}
% \setlength\evensidemargin{0mm}
% \setlength\topmargin{0mm}
% \setlength\headsep{0mm}
% \setlength\headheight{0mm}
% \let\openright=\cleardoublepage

%% Character encoding: usually latin2, cp1250 or utf8:
\usepackage[utf8]{inputenc}

%% Further useful packages (included in most LaTeX distributions)
\usepackage{amsmath}        % extensions for typesetting of math
\usepackage{amsfonts}       % math fonts
\usepackage{amsthm}         % theorems, definitions, etc.
\usepackage{bbding}         % various symbols (squares, asterisks, scissors, ...)
\usepackage{bm}             % boldface symbols (\bm)
\usepackage{graphicx}       % embedding of pictures
\usepackage{fancyvrb}       % improved verbatim environment
\usepackage[numbers]{natbib}         % citation style AUTHOR (YEAR), or AUTHOR [NUMBER]
\usepackage[nottoc]{tocbibind} % makes sure that bibliography and the lists
			    % of figures/tables are included in the table
			    % of contents
\usepackage{dcolumn}        % improved alignment of table columns
\usepackage{booktabs}       % improved horizontal lines in tables
\usepackage{paralist}       % improved enumerate and itemize
\usepackage[usenames]{xcolor}  % typesetting in color

\usepackage{multirow}
%\usepackage{obo-cite}

%%% Basic information on the thesis

% Thesis title in English (exactly as in the formal assignment)
\def\ThesisTitle{Automatic Error Correction of Machine Translation Output}

% Author of the thesis
\def\ThesisAuthor{Dušan Variš}

% Year when the thesis is submitted
\def\YearSubmitted{2016}

% Name of the department or institute, where the work was officially assigned
% (according to the Organizational Structure of MFF UK in English,
% or a full name of a department outside MFF)
\def\Department{Institute of Formal and Applied Linguistics}

% Is it a department (katedra), or an institute (ústav)?
\def\DeptType{Institute}

% Thesis supervisor: name, surname and titles
\def\Supervisor{RNDr. Ondřej Bojar, Ph.D.}

% Supervisor's department (again according to Organizational structure of MFF)
\def\SupervisorsDepartment{Institute of Formal and Applied Linguistics}

% Study programme and specialization
\def\StudyProgramme{Master of Computer Science}
\def\StudyBranch{Mathematical Linguistics}

% An optional dedication: you can thank whomever you wish (your supervisor,
% consultant, a person who lent the software, etc.)
\def\Dedication{%
Dedication.
}

% Abstract (recommended length around 80-200 words; this is not a copy of your thesis assignment!)
\def\Abstract{%
Abstract.
}

% 3 to 5 keywords (recommended), each enclosed in curly braces
\def\Keywords{%
{automatic post-editing,} {machine translation,} {supervised machine learning,} {natural language processing,} {Treex}
}

%% The hyperref package for clickable links in PDF and also for storing
%% metadata to PDF (including the table of contents).
\usepackage[pdftex,unicode]{hyperref}   % Must follow all other packages
\hypersetup{breaklinks=true}
\hypersetup{pdftitle={\ThesisTitle}}
\hypersetup{pdfauthor={\ThesisAuthor}}
\hypersetup{pdfkeywords=\Keywords}
\hypersetup{urlcolor=blue}


%%%% Custom Definitions %%%

\newcommand{\fixme}[1]{\textcolor{red}{FIXME: (#1)}} % macro for fixme entries
\newcommand{\todo}[1]{\textcolor{blue}{TODO: (#1)}} % macro for todo entries

\def\samp#1{``\textit{#1}''}
\def\pojem#1{\textit{#1}}

\def\Tref#1{Table~\ref{#1}}
\def\Fref#1{Figure~\ref{#1}}
\def\Eref#1{Example~\ref{#1}}
\def\Sref#1{Section~\ref{#1}}
\def\Cref#1{Chapter~\ref{#1}}
\def\equo#1{``#1''}
\def\notion#1{{\emph{#1}}}
\def\perscite#1{\newcite{#1}}
\def\parcite#1{\cite{#1}}
\def\footurl#1{\footnote{\url{#1}}}
\def\hash{\#}
\def\tilda{\~{}}


%%%%%%%%%%%%%%%%%%%%%%%%%%%

% Definitions of macros (see description inside)
%%% This file contains definitions of various useful macros and environments %%%
%%% Please add more macros here instead of cluttering other files with them. %%%

%%% Minor tweaks of style

% These macros employ a little dirty trick to convince LaTeX to typeset
% chapter headings sanely, without lots of empty space above them.
% Feel free to ignore.
\makeatletter
\def\@makechapterhead#1{
  {\parindent \z@ \raggedright \normalfont
   \Huge\bfseries \thechapter. #1
   \par\nobreak
   \vskip 20\p@
}}
\def\@makeschapterhead#1{
  {\parindent \z@ \raggedright \normalfont
   \Huge\bfseries #1
   \par\nobreak
   \vskip 20\p@
}}
\makeatother

% This macro defines a chapter, which is not numbered, but is included
% in the table of contents.
\def\chapwithtoc#1{
\chapter*{#1}
\addcontentsline{toc}{chapter}{#1}
}

% Draw black "slugs" whenever a line overflows, so that we can spot it easily.
\overfullrule=1mm

%%% Macros for definitions, theorems, claims, examples, ... (requires amsthm package)

\theoremstyle{plain}
\newtheorem{thm}{Theorem}
\newtheorem{lemma}[thm]{Lemma}
\newtheorem{claim}[thm]{Claim}

\theoremstyle{plain}
\newtheorem{defn}{Definition}

\theoremstyle{remark}
\newtheorem*{cor}{Corollary}
\newtheorem*{rem}{Remark}
\newtheorem*{example}{Example}

%%% An environment for proofs

%%% FIXME %%% \newenvironment{proof}{
%%% FIXME %%%   \par\medskip\noindent
%%% FIXME %%%   \textit{Proof}.
%%% FIXME %%% }{
%%% FIXME %%% \newline
%%% FIXME %%% \rightline{$\square$}  % or \SquareCastShadowBottomRight from bbding package
%%% FIXME %%% }

%%% An environment for typesetting of program code and input/output
%%% of programs. (Requires the fancyvrb package -- fancy verbatim.)

\DefineVerbatimEnvironment{code}{Verbatim}{fontsize=\small, frame=single}

%%% The field of all real and natural numbers
\newcommand{\R}{\mathbb{R}}
\newcommand{\N}{\mathbb{N}}

%%% Useful operators for statistics and probability
\DeclareMathOperator{\pr}{\textsf{P}}
\DeclareMathOperator{\E}{\textsf{E}\,}
\DeclareMathOperator{\var}{\textrm{var}}
\DeclareMathOperator{\sd}{\textrm{sd}}

%%% Transposition of a vector/matrix
\newcommand{\T}[1]{#1^\top}

%%% Various math goodies
\newcommand{\goto}{\rightarrow}
\newcommand{\gotop}{\stackrel{P}{\longrightarrow}}
\newcommand{\maon}[1]{o(n^{#1})}
\newcommand{\abs}[1]{\left|{#1}\right|}
\newcommand{\dint}{\int_0^\tau\!\!\int_0^\tau}
\newcommand{\isqr}[1]{\frac{1}{\sqrt{#1}}}

%%% Various table goodies
\newcommand{\pulrad}[1]{\raisebox{1.5ex}[0pt]{#1}}
\newcommand{\mc}[1]{\multicolumn{1}{c}{#1}}

%%% Example environment
\newcounter{excounter}
    \newenvironment{myexample}[0]{%
    \bigskip\noindent%
    \refstepcounter{excounter}%
    \newline%
}
{%
    \begin{center}%
    \textsc{\theexcounter}%
    \end{center}%
    \par\bigskip%
}
\numberwithin{excounter}{chapter}


% Title page and various mandatory informational pages
\begin{document}
%%% Title page of the thesis and other mandatory pages

%%% Title page of the thesis

\pagestyle{empty}
\hypersetup{pageanchor=false}
\begin{center}

\centerline{\mbox{\includegraphics[width=166mm]{img/logo.pdf}}}

\vspace{-8mm}
\vfill

{\bf\Large MASTER THESIS}

\vfill

{\LARGE\ThesisAuthor}

\vspace{15mm}

\begin{center}
\LARGE\bfseries\ThesisTitle
\end{center}
%{\LARGE\bfseries\ThesisTitle}

\vfill

\Department

\vfill

\begin{tabular}{rl}

Supervisor of the master thesis: & \Supervisor \\
\noalign{\vspace{2mm}}
Study programme: & \StudyProgramme \\
\noalign{\vspace{2mm}}
Study branch: & \StudyBranch \\
\end{tabular}

\vfill

% Zde doplňte rok
Prague \YearSubmitted

\end{center}

\newpage

%%% Here should be a bound sheet included -- a signed copy of the "master
%%% thesis assignment". This assignment is NOT a part of the electronic
%%% version of the thesis. DO NOT SCAN.

%%% A page with a solemn declaration to the master thesis

\openright
\hypersetup{pageanchor=true}
\pagestyle{plain}
\pagenumbering{roman}
\vglue 0pt plus 1fill

\noindent
I declare that I carried out this master thesis independently, and only with the cited
sources, literature and other professional sources.

\medskip\noindent
I understand that my work relates to the rights and obligations under the Act No.~121/2000 Sb.,
the Copyright Act, as amended, in particular the fact that the Charles
University has the right to conclude a license agreement on the use of this
work as a school work pursuant to Section 60 subsection 1 of the Copyright Act.

\vspace{10mm}

\hbox{\hbox to 0.5\hsize{%
In ........ date ............	% FIXME!
\hss}\hbox to 0.5\hsize{%
signature of the author
\hss}}

\vspace{20mm}
\newpage

%%% Mandatory information page of the thesis

\openright

\vbox to 0.5\vsize{
\setlength\parindent{0mm}
\setlength\parskip{5mm}

Title:
\ThesisTitle

Author:
\ThesisAuthor

\DeptType:
\Department

Supervisor:
\Supervisor, \SupervisorsDepartment

Abstract:
\Abstract

Keywords:
\Keywords

\vss}

\newpage

%%% Dedication

\openright

\noindent
\Dedication

\newpage

\openright
\pagestyle{plain}
\pagenumbering{arabic}
\setcounter{page}{1}


%%% A page with automatically generated table of contents of the master thesis

\setcounter{page}{1}
\tableofcontents

%%% Each chapter is kept in a separate file
\chapter{Introduction}
%\chapter*{Introduction}
%\addcontentsline{toc}{chapter}{Introduction}

In this thesis, we are presenting a statistical post-editing tool MLFix
based on its rule-based predecessor Depfix \citep{depfix:2014}.
We aim to use statistical machine learning methods to generalize a subset of Depfix rules
with focus on creating fairly language-independent automatic post-editing (APE) tool.
Our main goal is to find a reasonable tradeoff between the amount
of linguistic knowledge gathered from the input data and the dependence
on the language-specific third-party analysis tools, to specify the
machine learning (ML) task which our APE component should accomplish and to train
a sufficient ML model for correcting the machine translation (MT) output.
The system was developed using English-Czech language pair, however, we also present
a preliminary results we gathered from English-German language pair experiments.

\section{Task Motivation}

Even though the current state-of-the-art MT systems have been gradually improving in the
past few decades, they are still not perfect. Currently the most popular statistical machine
translation (SMT) systems can be fairly effective even without any or a little initial
linguistic knowledge about the concerned language pair given they have access to sufficient
amount of parallel data. However, when translating into the morphologically rich languages,
the data sparsity increases rapidly and such systems quickly begin to introduce grammatical
errors into the translated sentences and worsen the overall fluency of the translation.

Still, these systems can be of help to the human translators since it is usually less costly
to have even partially incorrect translations provided by the SMT system and have
human translator correct them then to translate the source text manually.
Human post-editors can however still be a quite costly so naturally there have been
attemts to automate this process with automatic post-editing tools.

\section{Related Work}

Up to this date there have been various approaches to the task of automatic
error correction of the machine translation output, each with a different level of success.

The very first attempts in the field of automatic post-editing \citep{simard2007rule}
focused on applying phrase-based statistical machine translation (PB-SMT)
system on the output of the rule-base machine translation. The system was trained on a monolingual
bitextual data containing the MT output as a source sentence and a reference translation as
the target sentence. This phrase-based automatic post-editor (PB-APE) helped to significantly improve
the performance of the rule-based MT system in question.
There were also attempts to apply the post-editing component
on a phrase-based translation system but the combined system performed slightly worse in comparison
with the standalone SMT system.

Further experiments with PB-SMT $+$ PB-APE \citep{bechara:2011} were done
on English-French and French-English translation, where significant improvements were reported
for the latter translation direction. They improved the design suggested by Simard et al. by
creating the purely statistical pipeline. The PB-APE was then expanded by adding additional
context information about the source sentence to SMT-generated output (used as an input for
the PB-APE) reporting further improvement in performance. However, the presented system failed to
improve the MT output in the opposite translation direction (English-French).

These previous attempts can be considered purely statistical since only a little ammount of linguistic
knowledge about the concerned languages was used during developement. To gain further insight into
the task of automatic post-editing, more thorough analysis of the most frequent errors made by the current
SMT systems was performed by \citet{bechara:master} and later by \citet{biblio:RoAutomaticpostediting2013}, the former
for English-French and the latter for English-Czech.

The error analysis was later used during developement of the rule-based APE Depfix \citep{depfix:2014},
which was designed to correct the errors made by the English-Czech SMT systems.
The system uses a set of finely hand crafterd rules that aim at identifying
and correcting morphological errors such as incorrect agreement or valency, which are
often encountered during English-Czech machine translation. It does not focus on a lexical errors
although some minor corrections, e.g. insertion of missing reflexive verbs, are made. The system
succeeded to improve output of various MT systems and was deployed as a stable part of the
Chimera \citep{biblio:BoRoChimera2013} MT system.

The idea behind Depfix system seems promising, however, due to its rule-based nature, it is difficult
to apply the APE on different language pair since it would required costly modifications to the
existing set of rules. One of the goals of this thesis is to try and replace the rule-based blocks
by a statistical ones so they applied to the different target language MT output more easily simply
by training an appropriate statistical model.

%"Better Evaluation for Grammatical Error Correction" Dahlmeier
% napsat az do model tuningu ???
% Correction model for non-MT generated errors (classification)
% Important -> balancing training data

\section{Thesis Structure}
%strucny popis depfixu? popis castych chyb ???

%odkazy
In~Chapter~\ref{chap:system_descr}, we describe MLFix data processing
pipeline and introduce all the tools we use.
In~Chapter~\ref{chap:data}, we take a look at all datasets which we used during our system
developement and analyze their usefulness for our system.
In~Chapter~\ref{chap:task_descr}, we define the post-editing task, go through
various modifications of the task we considered and explain our though
process with support of the analysis of the available data.
In~Chapter~\ref{chap:tuning}, we describe the process of the developement
the statistical models for the MLFix post-editing component and present
a preliminary results of the evaluation of the trained models.
In~Chapter~\ref{chap:eval}, we evaluate the performance of the whole MLFix system
and analyze the level of contribution to the resulting performance
of each individual statistical compocomponent.
In~Chapter~\ref{chap:german}, we briefly describe our modifications to English-Czech pipeline
when used on English-German language pair, then summarize the differences
in the training data and model training between the two language pairs. The Chapter concludes with evaluation of English-German pipeline.
%In chapter 8, we present performance results of the modified English-German
%pipeline.
We conclude the thesis in Chapter 9.

We include a full English-Czech and English-German Treex scenario in Attachment~\ref{attach:scen}.

\fixme{Dodatek k attachmentum}



\chapter{System description}
\label{chap:system_descr}

In this chapter we describe the main components
of MLFix system and take a closer look at the suggested processing pipeline.
We focus on English-Czech pipeline, however, we also describe modifications
needed to apply MLFix for other language pairs.

Full examples of currently deployed Treex scenarios (English-Czech and English-German)
are included in Attachment~\ref{attach:scen}.

%The main idea of MLFix system is to analyze the provided input data providing
%a set of input features 

\section{Processing pipeline}

MLFix is, similarly to its predecessor, almost entirely implemented in the
Treex \citep{Popel:2010:TMN:1884371.1884406}\footurl{http://ufal.mff.cuni.cz/treex}
framework (formerly knownas TectoMT).
The framework was originally created as English-Czech hybrid translation system, combining
rule-based modules with statistical models and using deep semantic language representation
for the sentence translation. However, due to its modularity, it is
now used for various tasks of natural language processing (NLP) across different
languages. The framework was built to support a methodology of the theory of Functional Generative Description \citep{Sgall1967}
and was adapted to support sentence representation in Prague Dependency Treebank \citep{pdt20:2006}.
Mainly, it supports the representation of sentences on different layers of abstaction defined in FGD: morphological layer,
analytical layer and a tectogrammatical layer.\footnote{Usually referred to as m-layer, a-layer and t-layer where prefixes m-, a- and t- are also used to refer to the objects on the corresponding layer of abstraction.}
Because the tools currently available for tectogrammatical layer analysis are available for
a limited set of languages (mainly Czech and English, a developement of others is in progress),
we decided to use only the a-layer (surface syntax) for data representation.

\subsection{M-Layer analysis}

MLFix analysis pipeline is derived from an existing Depfix pipeline
with a several modifications to make it easier to apply to different
target languages. Input data (source sentence + MT output aligned on a sentence level,
or additionaly reference sentences for extracting the training data)
are first read in parallel and stored into Treex internal representation.
Both source side and MT side are tokenized by a rule-based tokenizer, each token is then
represented by a separate m-node.

Next, a lemmatization and part-of-speech (POS) tagging is performed on each m-node.
For both English and Czech, we use MorphoDiTa \citep{strakova14:morphodita}\footurl{http://ufal.mff.cuni.cz/morphodita}
tool for morphological analysis and tagging. MorphoDiTa is also used for Czech lemmatization.
For English lemmatization, we use a rule-based block implemented in Treex.
It is important to provide a POS tagger for the target language that supports
relatively fine-grained morphological tags because our goal is to correct morphological
errors represented mainly through these tags.
It was reported by Rosa\fixme{znovu citovat diplomku?} that tagger produces significantly more
errors in morphological analysis of Czech SMT outputs. In Depfix, this is covered
by a rule-based block that identifies these errors and changes the incorrect POS tag without
changing the surface form. We decided to omit this block and leave the issue to our statistical
component.

Last step we do in the scope of the m-layer analysis is a transformation of the POS tags into
a more general representation. Optionally we can apply a named-entity recognition
tool if one is available but it isn't mandatory. For English, we use the Stanford
Named Entity Recognizer (NER) \citep{Finkel:2005:INI:1219840.1219885}\footurl{http://nlp.stanford.edu/software/CRF-NER.shtml},
for Czech, we use a simple rule-based NER. 

\subsection{Interset}

Since there are usually different tagsets used across individual languages,
often engineered for purposes of that specific language and with no standardized
tag representation,
we would be forced to modify our existing pipeline to some extent everytime
a new language would be introduced.
Therefore we decided to use Interset \citep{biblio:ZeReusableTagset2008}\footurl{https://ufal.mff.cuni.cz/interset}, an interlingua-based
representation of POS tags from various tagsets. To be able to use this
representation to represent tags from a given tagset, a decoding/encoding
module is required. However, the support for various tagsets spanning
through different languages started growing lately mainly due to Universal Dependencies \citep{universal-dep:2016}\footurl{universaldependencies.org} project.

After the transformation, in following steps of the analysis, a choice
of a specific tagset becomes transparent because MLFix blocks only
have to deal with one well-defined set of features.

\subsection{Word alignment}

In the next step, we create a word-level alignment between each sentence pair
using GIZA++ \citep{och:ney:2000}. We make one-to-one word alignment where possible
using an intersection symmetrization. This step helps us later with feature extraction
and with further processing of the target sentence.

In the process of training data extraction, we also create a simple alignment between
the MT sentence and the reference sentence exploiting forms, lemmas and tags
of the m-nodes. The alignment between the source sentence and the MT ouptut is then also
projected to the reference sentence.

\subsection{A-Layer analysis}

After the m-layer analysis and the word alignment, we perform a dependency parsing.
For English, we use Maximum spanning tree (MST) parser \citep{mcdonald:pereira:ribarov:hajic:2005}\footurl{http://sourceforge.net/projects/mstparser}
implemented in Treex framework. For SMT output, even though there might
be an existing dependency parser available for the target language, it is usually
trained on data that do not contain errors. Therefore, it has usually significantly
lower performance when applied on the SMT output. Still research on Depfix have shown
that a knowledge of the dependency structure of the SMT output can provide additional
valuable information for identifying grammatical errors\footnote{Actually, in the case of Depfix
, the information about the dependency structure is crucial for most of the fixing components
because e.g. a parent-child relationship is examined almost everytime.}, thus improving
the post-editor's performance.

For the time being we have decided to build the dependency structures of the SMT output simply
by projecting the dependency structure of the source sentence onto the target side
using the word alignment we extracted in the preceeding step. The resulting structure
is likely to contain errors, but should be at least consistent thoughout our data.
To compensate for the lower accuracy of the extracted dependency structure we perform
a dependency parsing of the reference sentences during training, if a proper parser
is available. For this purposes, we use the MST Parser for the Czech language as well.
We expect that the combination of dependency parser for the reference sentences and
a right choice of constraints applied when extracting training instances we should
avoid polluting our training data with instances containing misleading context information
due to an incorrect structure of the SMT dependency tree.

For Czech SMT output, an implementation of the MST parser adapted for the SMT output
is already available \citep{biblio:RoDuUsingParallel2012}, however, so far we haven't done any experiments
to determine whether the dependency structures provided by the adapted parser (which
should be more accurate) influence the final performance of our system.

%\subsection{Training pipeline}

%We also do preprocessing of data we use for training our statistical component.
%The pipeline is similar to the processing pipeline with addition of processing
%the reference sentences. The reference sentences are processed in a same way
%as the MT output, however they are not analyzed on the a-layer, only a simple
%lemma based word alignment with the MT output is made.

\section{Statistical component}

After data preprocessing, we extract all available features and
apply a trained statistical model. The features are extracted separetely for each node and passed
to a model. The model needs to accomplish these two goals:
\begin{enumerate}
    \item identify candidates with incorrect morphology,
    \item fix incorrect morphological features.
\end{enumerate}
Of course, these two steps can be split between multiple separate components (and models).
We describe the post-editing process in more detail later in chapter~\ref{chap:tuning}.

We decided to use Scikit-Learn \citep{scikit-learn}\footurl{biblio:RoDuUsingParallel2012} toolkit to train and execute our models since
it has easy-to-use and quite uniform interface, which allows us to try out different
machine learning (ML) methods simply by switching the model class. In Treex, we use a simple
wrapper to load and execute the trained Scikit-Learn model. If a support for a different ML implementation
is needed (e.g. VowpalWabbit\footurl{https://github.com/JohnLangford/vowpal\_wabbit/wiki}\footnote{At the time of writing this thesis, there is already
a limited support for VowpaWabbit.}), this wrapper can be easily modified to suit the requirements.

\section{Wordform generation}

After we have identified morphologically incorrect words and assigned
them a new morphological categories we need to generate a new surface form reflecting
the changes we have made.
This can be done either by rule-based component or with
help of another statistical model. 

For Czech, we have used a morphological generator build upon the morphology of \citet{HajicHAB2004}
that is already a part of Treex framework. For other languages,
when there was not another option available,
we have used Flect \citep{DBLP:conf/acl/DusekJ13}\footurl{https://ufal.mff.cuni.cz/flect}.
It is a language independent morphological generation tool also using Scikit-Learn
models. The tool learns morphological inflection patters form an annotated corpora and
should able to inflect even previously unseen words using lemma suffixes as features
and predicting a difference between a lemma and a surface form.

\section{Language independence}

From the previous description, we can notice that MLFix still depends on several
language specific tools. It is quite dependent on the capability of a source
language analysis (in our case English), because we do not only need a POS tagger but also a dependency
parser. However, we assume that the source sentences are usually grammatical correct
so it is much easier to provide required tools than their modified versions
targeted at the MT output.

We also require a specific decoder of the source and target POS tags into Interset feature
structures but Interset already covers large variety of the most widespread tagsets available
and its support is still growing.

Finaly, aside from a language specific post-editing models (which we do not expect
to be reusable across different languages) we require a module that regenerates
a correct wordform (either from the \samp{lemma+tag} or \samp{lemma+Interset features combination}).
A state-of-the-art tool might not be explicitly available for every language but if no
other option is provided we can use a statistical form generator, in our case Flect.

\chapter{Available data}

% dostupna data
% mnozstvi
% popis
% vyuziti

XXX TOTO ZREJME AZ DO POPISU SYSTEMU

In this chapter we are going to describe the format of the input data, that
our system require for training and present all the available datasets.
The differences between the suggested data can seem minor but they can
have impact on the overall performance of the system.

To train our system, we require sentence-level aligned set of source sentences,
MT output and a reference sentences
%TODO: footnote - reference in this thesis as a triparallel data
. The reference sentences can be either
a classic reference translation of the source text or a result of a human
post-editation of the MT output, they are however required because they
provide use with the possible corrections of the MT output.
Naturally, the closer the MT output is
to the reference translation, the easier it should be for our system
to extract valuable learning instances form the sentences. We also
considered using only pair of sentence aligned pair of MT translated
sentences and the reference sentences however, this way the trained
model will lose some of the useful features which can be extracted from
the triparallel data.

XXX

In this chapter we are going to take a closer look at the available
sources of data and describe how they contributed to our research.

We came across a various sources of training data with various level
of usefulness and usually of a smaller volume. Some of the sources
are:
\begin{itemize}
\item Khan's school human
post-edits of manually translated (EN$\rightarrow$CS) subtitles,
\item the Autodesk
triparallel data\footurl{https://autodesk.app.box.com/Autodesk-PostEditing},
\item log files of human post-editing done by Lingea for the
HimL\footurl{http://www.himl.eu/} project test dataset,

% uvest QT21?
%\item data from the QT21 project\footurl{http://www.qt21.eu/deliverables/annotations/},

%TODO: wmt - odkaz, footnote?
\item results from the previous workshops on machine translation (WMTs).
\end{itemize}

In the following sections, we will describe in more detail each of the datasets.

\subsection{Khan's school}

The data provided by Khan's school consist of en-cs subtitles,
where the Czech part (usually manually translated from English) was manually
edited. During the analysis of the dataset,
we've noticed that most of the time, the corrections were made
mostly on a lexical level which is only natural since the Czech sentences
were created by a human translator.
Therefore, we decided to avoid using this dataset for the time being or to
rather treat the corpus as a simple bilingual data.

\subsection{Autodesk}

Autodesk data consist of English sentences which were machine translated into
a set of target languages (cs, de, pl etc.) complemented with human post-editing
of MT output. These datasets are domain specific (mostly user documentations)
however, so they might not be very attractive to use with more general texts.
We weren't able to gather any information about the MT system that was used
to create the translated output. The biggest advantage of these data is
their larger volume when compared with other post-editing datasets so we
used them mainly for model developement and benchmarking of the used machine
learning methods.

\subsection{HimL-Lingea logs}

The data provided by Lingea were collected when official HimL test sets were
created. The data consist mainly of the public health related texts.
The original English sentences were first
machine-translated to HimL languages (Czech, German, Polish and Romanian)
and then post-edited by professional
translators using Lingea's post-editing tool. The datasets are probably the most
detailed one since they consist of complete logfiles
describing elementary actions taken by human post-editors (such as selecting
phrases in a translated sentence, looking up alternative translations
in a dictonary etc.).

When we examined the data more closely,
we noticed that it is rather difficult to determine which actions are
useful for our machine learning process. Also, we were little disapointed
when we found out that most of the time, the post-editors preferred to
rewrite the whole sentence "from scratch"\footnote{By that, we mean that the
translator usually preffered to write the whole corrected sentence, even though
only a several changes (both lexical and morphological), and delete the original.
This was also probably motivated by the need of reordering in the MT translated
sentence.} opposite to doing more atomic modifications to the provided MT output.

In the end, we decided to simply extract a triparallel data from these logs (the source
sentence + SMT output + result of the human post-editing). In the future,
we might consider to use other logged actions for model training.

\subsection{WMT datasets}

For the last decade, the workshop on machine translation (WMT) has aimed
to provide working grounds for many researchers in the field of machine
translation. It has been great source of the parallel data between English
on one side and various other languages on the other. Each year, the scope
of the workshop expands, including various new subtask related to machine translation,
such as several evaluation tasks and, more recently, automatic post-editing task.

The data available for the post-editing task usually contains a set of:
\begin{itemize}
    \item source English sentences,
    \item output of various MT systems, usually the ones that participate in the main
translation tasks,
    \item human post-edition of the mixed MT output from the participating MT systems.
    \fixme{je tohle pravda?}
\end{itemize}

These datasets gives us the opportunity to compare the performace of our post-editing
models when applied to the different systems. Even though the data available for the
WMT subtasks are from various domains (news, IT, biomedical), the domain of the post-editing
data is more limited, mainly to the news articles.

\subsection{Other sources}

The sources listed above (each one to a different degree, maybe except for the last one) can
be considered a knowledge base for examining the behaviour of a human post-editor as well
as training data for our system. We think that they might provide us
some interesting insight into the post-editor's thought process. On the other hand,
might also consider using some other sources since the sources mentioned above are quite
limited.

One possible way to face the data sparsity is to use available parallel corpora.
These corpora (containing only \notion{source sentences + reference translations})
can be expanded by translating the source sentences and thus creating a set of 
sentences which contain MT generated errors and should be fixed to resemble the 
reference translation. These data can be then used to train models for that
specific SMT system\footnote{Obviously, the post-editing model training data have to be
different from the parallel data used for the training of the SMT system, so
some
jackknife sampling should be used with limited training data.} or possibly for
another SMT systems.
This method can surely help to overcome the aforementioned data acquisition bottleneck
since there is generally much more parallel data then post-edited sentences. For
English-Czech language pair, the natural choice of the parallel corpus would be
CzEng~1.0\footurl{http://ufal.mff.cuni.cz/czeng} \cite{czeng10:lrec2012}.
We didn't use this approach in the scope of this thesis however, it might
be considered in our future work.
\todo{posledni vetu az do shrnuti?}

Of course, this aproach introduces some additional noise
related to the post-editing task. For example, we can get fluent MT output which
is just a variation on the reference translation (possibly thanks to the different
choice of wording etc.). It can be therefore hard to distinguish these fairly correct training
instances from the incorrect ones.

The basic summary of the available data is shown in the table \Tref{avail-data}.

\subsection{Monolingual data}

We also considered simply using bitext monolingual data (either \notion{MT output + post-edited sentences}
or \notion{MT output + reference translations}), however due to the nature of our processing
pipeline we would be much more limited when analyzing the training data. Also, as
mentioned by Rosa\cite{depfix:2014}, the additional information that can be
extracted from the source
sentences proved to be valuable in the practice.

%Even if we do not use this approach for the training of the final MLFix models,
%the volume of the data might help us with feature selection and preliminary model evaluation.

\begin{table*}[t]
\centering
\small

\begin{tabular}{lcccc}
\multirow{2}{*}{}  &  \multirow{2}{*}{\hash{} Sentences}  &  \multicolumn{3}{c}{\hash{} Tokens}  \\
&   & English & Czech (MT) & Czech (PE) \\
\hline
Khan's school & NA & \tilda{}93k & \tilda{}93k & \tilda{}93k \\
Autodesk & 46,916 & 490,005 & 456,697 & 441,645 \\
HimL-Lingea & 3892 & 60,142 & 51,428 & 56,485 \\
CzEng 1.0 & 15M & 206M & NA & 150M \\
\end{tabular}
\caption{Summary of the available post-editing data.
}
\label{avail-data}
\end{table*}

%The
%number of parallel sentences and the number of tokens in the English source, the
%MT output (\equo{MT}) and the post-edited MT output (\equo{PE}). Only
%English-Czech data is listed since these datasets for other target languages
%(where available) are similar in volume.
%For Khan's school, we only
%provide estimates. The CzEng data set is not translated by any SMT at the moment,
%so the related information is omitted. We provide the information only for comparison.

%We do not include information about the QT21 data since we are yet to explore them.



\chapter*{Conclusion}
\addcontentsline{toc}{chapter}{Conclusion}


%%% Bibliography
%\include{bibliography}
\bibliographystyle{plainnat}
%\bibliographystyle{obo-bst}
\phantomsection
\addcontentsline{toc}{chapter}{Literature}
\bibliography{bib-varis}


%%% Figures used in the thesis (consider if this is needed)
\listoffigures

%%% Tables used in the thesis (consider if this is needed)
%%% In mathematical theses, it could be better to move the list of tables to the beginning of the thesis.
\listoftables

%%% Abbreviations used in the thesis, if any, including their explanation
%%% In mathematical theses, it could be better to move the list of abbreviations to the beginning of the thesis.
\chapwithtoc{List of Abbreviations}

%%% Attachments to the master thesis, if any. Each attachment must be
%%% referred to at least once from the text of the thesis. Attachments
%%% are numbered.
%%%
%%% The printed version should preferably contain attachments, which can be
%%% read (additional tables and charts, supplementary text, examples of
%%% program output, etc.). The electronic version is more suited for attachments
%%% which will likely be used in an electronic form rather than read (program
%%% source code, data files, interactive charts, etc.). Electronic attachments
%%% should be uploaded to SIS and optionally also included in the thesis on a~CD/DVD.
\chapwithtoc{Attachments}

\openright
\end{document}
